% =====================================================================
%  Lecture 23 -- User Plane                               (English edition)
% =====================================================================
\documentclass[11pt,a4paper]{article}
\providecommand\ACNroot{../..}
\input{\ACNroot/templates/lecture-en}

\ACNset{number}{23}
\ACNset{title}{The 5G User Plane}
\ACNset{shorttitle}{User plane}
\ACNset{titleother}{صفحهٔ کاربر نسل پنجم}
\ACNset{part}{Part V -- Planes and outlook (Lectures 23--26)}
\ACNset{time}{6--8 hours (3 h reading, 3--5 h exercises)}
\ACNset{level}{Advanced}
\ACNset{prereq}{Lectures 6, 13, 16, 19 and 22}
\ACNset{gen}{5}

\begin{document}
\ACNcover

% ---------------------------------------------------------------------
\section{Learning objectives}
After studying this lecture you should be able to:
\begin{objectives}
  \item \textbf{Draw} the 5G user-plane protocol stack from UE to data network, including the gNB's CU/DU split.
  \item \textbf{Explain} the roles of SDAP, PDCP, RLC and MAC, and of GTP-U on F1-U, N3 and N9.
  \item \textbf{Compare} RAN functional splits and \textbf{estimate} their fronthaul requirements.
  \item \textbf{Describe} UPF packet processing with PFCP rules and its implementation options.
  \item \textbf{Compute} header overhead, latency budgets, duplication reliability and sequence-number windows.
\end{objectives}

% ---------------------------------------------------------------------
\section{Prerequisite knowledge}
You need the LTE radio protocol stack (Lecture~6), GTP-U and PFCP (Lecture~13), NFV data-plane techniques (Lecture~16), UPF topologies
(Lecture~19) and the QoS model (Lecture~22). This lecture and the next one look at the 5G system ``by plane'': first the path of user
data, then the signalling that controls it.

\paragraph{Key terms} The vocabulary of this lecture with its Persian equivalents:
\begin{keyterms}
User plane & \fa{صفحهٔ کاربر} & Functions and protocols that carry user data. \\
CU / DU / RU & \fa{واحد مرکزی / توزیع‌شده / رادیویی} & Parts of a split gNB. \\
F1-U / N3 / N9 & \fa{واسط‌های صفحهٔ کاربر} & GTP-U interfaces CU--DU, RAN--UPF, UPF--UPF. \\
Functional split & \fa{تقسیم کارکردی} & Point in the protocol stack where the RAN is divided. \\
Fronthaul & \fa{فرانت‌هال} & Transport between DU and RU (or baseband and radio). \\
PDCP duplication & \fa{تکرار PDCP} & Sending the same PDCP PDU over two legs for reliability. \\
SN window & \fa{پنجرهٔ شمارهٔ ترتیب} & Range of sequence numbers in flight for reordering. \\
PDR / FAR / QER / URR & \fa{قواعد PFCP} & Detection, forwarding, QoS and usage rules in the UPF. \\
DPDK / SmartNIC & \fa{شتاب صفحهٔ داده} & Software and hardware acceleration of packet processing. \\
Local breakout & \fa{خروج محلی} & Exiting to a data network close to the RAN. \\
\end{keyterms}

% ---------------------------------------------------------------------
\section{Introduction and motivation}
Everything else in a mobile network exists to serve the user plane: the path that carries the subscriber's packets. Its performance
-- throughput, latency, reliability, cost per bit -- is what users experience and what operators pay for. In 5G the user plane is more
flexible than ever: the gNB can be split into centralised and distributed units, the core user plane is a programmable UPF that can be
placed anywhere from a central data centre to a factory, and several UPFs can be chained in one session.

This flexibility creates design questions that this lecture equips you to answer quantitatively: where to split the RAN given the
available fronthaul, where to place UPFs given a latency budget, how to reach URLLC reliability with duplication, and how much the protocol
headers cost \cite{ts38300,ts38401,ts23501}.

% ---------------------------------------------------------------------
\section{Fundamental concepts}

\subsection{The end-to-end stack}
User data pass (\cref{fig:stack}): in the UE through SDAP, PDCP, RLC, MAC and PHY; over the radio to the gNB-DU (RLC, MAC, PHY); over
F1-U in GTP-U to the gNB-CU-UP (PDCP, SDAP); over N3 in GTP-U to the UPF, possibly over N9 to further UPFs; and out of N6 to the data
network. The radio layers have these roles \cite{ts38300}:
\begin{itemize}
  \item \textbf{SDAP}: maps QoS flows to data radio bearers and marks QFIs (Lecture~22).
  \item \textbf{PDCP}: header compression (ROHC), ciphering and optional integrity protection, reordering and duplicate removal, PDCP
        duplication.
  \item \textbf{RLC}: segmentation and, in acknowledged mode, ARQ retransmission; NR's RLC no longer concatenates or reorders.
  \item \textbf{MAC}: multiplexing of logical channels, HARQ, scheduling.
\end{itemize}

\begin{acnfigure}
  \resizebox{\linewidth}{!}{\acnfig{lecture-23/up-stack}}
  \caption{5G user-plane protocol stack with a split gNB (CU-UP and DU).}
  \label{fig:stack}
\end{acnfigure}

\subsection{Splitting the gNB}
3GPP standardised a \emph{higher-layer split} between a central unit (CU: RRC, SDAP, PDCP) and a distributed unit (DU: RLC, MAC, PHY),
connected by F1 \cite{ts38401}; the CU is further split into control (CU-CP) and user plane (CU-UP), connected by E1. Lower-layer splits
between the DU and a radio unit (RU) -- in particular option 7.2, specified by the O-RAN Alliance -- move part of the physical layer into
the RU (\cref{fig:splits}). The lower the split, the more functions can be centralised (simpler sites, coordination across cells), but the
higher the fronthaul rate and the tighter the latency requirement (\cref{sec:fh}).

\begin{acnfigure}
  \resizebox{\linewidth}{!}{\acnfig{lecture-23/splits}}
  \caption{Selected RAN functional splits with indicative fronthaul rates (assumptions in \cref{sec:fh}).}
  \label{fig:splits}
\end{acnfigure}

\subsection{The UPF}
The UPF processes every packet against the rules installed by the SMF (\cref{fig:upf}) \cite{ts29244}: a \emph{PDR} detects the packet
(by TEID, UE address, SDF filter or QFI; the matching PDR with the highest precedence wins) and points to a \emph{FAR} (forward with a new
outer header, buffer, drop, duplicate), \emph{QERs} (gating and rate limits) and \emph{URRs} (counting and reporting volumes, time and
events). Buffering for idle UEs follows a \emph{BAR}. The UPF thus combines the S-GW and P-GW user-plane functions of the EPC (Lecture~13).

\begin{acnfigure}
  \resizebox{\linewidth}{!}{\acnfig{lecture-23/upf-pipeline}}
  \caption{UPF packet processing with PFCP rules.}
  \label{fig:upf}
\end{acnfigure}

% ---------------------------------------------------------------------
\section{Detailed technical explanation}

\subsection{PDCP duplication and multi-connectivity}
For URLLC, PDCP can duplicate packets and send the copies over two independent legs: two carriers of the same gNB (carrier aggregation)
or two gNBs (dual connectivity). The receiver's PDCP discards the second copy. Duplication converts the independence of the legs into
reliability (\cref{sec:dup}) and removes the interruption of a handover, at the cost of doubling the radio resources used for the duplicated
flow. It can be activated dynamically by MAC control elements when the channel deteriorates.

\subsection{Sequence numbers and reordering}
PDCP numbers each SDU (12- or 18-bit SN); RLC in acknowledged mode numbers segments (12- or 18-bit SN), in unacknowledged mode 6- or 12-bit. The
receiving PDCP delivers in order and waits up to a reordering timer for missing packets. The window of SNs that may be in flight is half the
SN space; at high data rates a short SN wraps quickly, which is why NR introduced 18-bit SNs (\cref{sec:sn}).

\subsection{F1-U and N3 transport}
On F1-U the CU-UP and DU exchange PDCP PDUs in GTP-U with an NR user-plane protocol header that carries flow-control information (e.g.\ the DU's
desired buffer size), so that the CU can pace the DU and avoid overflows when radio conditions change. On N3 and N9 the GTP-U header carries
the QFI and the reflective QoS indicator in an extension header. Both run over IP transport networks that must preserve QoS by mapping
QFIs/5QIs to DSCP values.

\subsection{UPF implementation and placement}
A UPF must classify, count and re-encapsulate every packet at line rate. Implementations range from software on servers with DPDK and SR-IOV
(Lecture~16), through SmartNIC offload of encapsulation and classification, to programmable switches running P4 pipelines (Lecture~15) for
the fast path while a CPU handles exceptions. Placement follows the latency budget (\cref{sec:budget}): central UPFs for Internet traffic,
regional or edge UPFs for latency-sensitive and locally served traffic, on-premises UPFs for enterprise and industrial networks.

% ---------------------------------------------------------------------
\section{Mathematical models}

\subsection{Protocol overhead}\label{sec:oh}
\paragraph{Model} An IP packet of $L$ bytes acquires SDAP, PDCP, RLC and MAC headers of $H = H_{\mathrm{SDAP}} + H_{\mathrm{PDCP}} + H_{\mathrm{RLC}} +
H_{\mathrm{MAC}}$ bytes on the radio. The efficiency is
\begin{equation}\label{eq:overhead}
  \eta = \frac{L}{L + H} .
\end{equation}
\paragraph{Example} SDAP 1, PDCP 3 (18-bit SN), RLC AM 3, MAC subheader 3 bytes: $H = 10$\,bytes. For $L = 1500$: $\eta = 99.3\,\%$; for a 60-byte voice
packet (after ROHC compresses its 60-byte IP/UDP/RTP header to a few bytes, $L$ is smaller still): $\eta = 85.7\,\%$. Integrity protection adds a
4-byte MAC-I to PDCP.

\subsection{Fronthaul rate}\label{sec:fh}
\paragraph{Model} With split option 8 the fronthaul carries time-domain IQ samples: sampling rate $f_s$, $b$ bits per I and per Q component, $N_{\mathrm{ant}}$
antennas and line coding $c$:
\begin{equation}\label{eq:cpri}
  R_{\mathrm{FH},8} = 2\, f_s\, b\, N_{\mathrm{ant}}\, c .
\end{equation}
With split 7.2 the fronthaul carries frequency-domain samples of the occupied resource elements per layer, approximately
$R_{\mathrm{FH},7.2} \approx 12\,N_{\mathrm{PRB}}\,N_{\mathrm{sym}}\,N_{\mathrm{layers}}\,2b'$ with $N_{\mathrm{sym}}$ symbols per second and compressed bit width $b'$.
\paragraph{Example} 100\,MHz NR at 30\,kHz: $f_s = 122.88$\,Msample/s, $b = 16$, 4 antennas, $c = 66/64$: $R_{\mathrm{FH},8} = 16.2$\,Gbit/s; with 64
antennas 260\,Gbit/s -- impractical. Split 7.2 with 273 PRBs, 28\,000 symbols/s, 4 layers, $b' = 9$: about 6.6\,Gbit/s, scaling with layers rather than
antennas. The higher-layer split (option 2) carries roughly the user data rate (about 2.5\,Gbit/s for the carrier of Lecture~14).
\paragraph{Interpretation} Massive MIMO is only practical with splits that do beamforming in the RU; the choice of split is a trade-off between
centralisation and transport cost.

\subsection{Latency budget}\label{sec:budget}
\paragraph{Model} One-way user-plane latency between UE and an application server is the sum of its components:
\begin{equation}\label{eq:budget}
  T = T_{\mathrm{radio}} + T_{\mathrm{gNB}} + T_{\mathrm{FH}} + T_{\mathrm{BH}} + T_{\mathrm{UPF}} + T_{\mathrm{N6}} .
\end{equation}
\paragraph{Example} With illustrative values (\cref{fig:budget}): radio 2\,ms, gNB processing 0.5\,ms, fronthaul 0.1\,ms, UPF 0.1\,ms; an edge UPF with 0.2\,ms
transport and 0.1\,ms to the server gives 3.0\,ms, a central UPF 600\,km away (3\,ms transport) with 1\,ms to the server gives 6.7\,ms. Only the edge
placement meets a 4\,ms one-way target.

\begin{acnfigure}
  \acnfig{lecture-23/latency-budget}
  \caption{One-way user-plane latency budget with edge and central UPF (illustrative values, \cref{eq:budget}).}
  \label{fig:budget}
\end{acnfigure}

\subsection{Reliability with duplication}\label{sec:dup}
\paragraph{Model} If each of two independent legs loses a packet with probability $p_1$, $p_2$, duplication loses it only if both fail:
\begin{equation}\label{eq:dup}
  P_{\mathrm{loss}} = p_1\,p_2 .
\end{equation}
\paragraph{Example} $p_1 = p_2 = 10^{-3}$ gives $10^{-6}$, meeting a URLLC target that a single leg could reach only with a much more robust (and costly)
MCS or more retransmissions -- which might not fit the delay budget.
\paragraph{Limitation} Correlated losses (both carriers from the same site blocked by the same obstacle) break the independence assumption; dual
connectivity with geographically separated nodes is more robust than carrier-level duplication.

\subsection{Sequence-number window}\label{sec:sn}
\paragraph{Model} With an $s$-bit SN, at most $2^{s-1}$ PDUs may be in flight. At a packet rate $P$ the window covers
\begin{equation}\label{eq:window}
  T_{\mathrm{win}} = \frac{2^{s-1}}{P} .
\end{equation}
\paragraph{Example} 2\,Gbit/s with 1500-byte packets: $P = 1.67\times10^5$\,packets/s. A 12-bit SN covers 12.3\,ms -- shorter than HARQ and ARQ delays plus
reordering; an 18-bit SN covers 0.79\,s. High-throughput bearers therefore use 18-bit SNs.

% ---------------------------------------------------------------------
\section{Network architecture and diagrams}
A typical 5G deployment combines RUs at the antenna sites, DUs at the cell sites or in nearby hubs (connected by 7.2 fronthaul), CU-UPs in regional
data centres together with edge UPFs, and central UPFs for Internet traffic. Each element can be virtualised on common infrastructure
(vRAN, Lecture~16), and the SMF chains UPFs per session (\cref{fig:stack,fig:splits,fig:upf}).

% ---------------------------------------------------------------------
\section{Step-by-step procedures}

\begin{procedurebox}{Following a downlink packet}\label{proc:dl}
\begin{steps}
  \item The packet arrives on N6 at the PSA UPF; a PDR matches it (UE address and SDF filter), the QER enforces rates, the URR counts it, and the FAR
        encapsulates it in GTP-U with the QFI towards the next hop (N9 or N3).
  \item The CU-UP receives it on N3, SDAP maps the QFI to a DRB, PDCP compresses headers, ciphers and numbers it, and sends it in GTP-U over F1-U.
  \item The DU's RLC segments it as scheduled by MAC; MAC multiplexes it into a transport block; PHY transmits it with HARQ.
  \item The UE's layers reverse the process; PDCP reorders and delivers in sequence to SDAP and the IP layer.
\end{steps}
\end{procedurebox}

\begin{procedurebox}{Choosing a RAN split}\label{proc:split}
\begin{steps}
  \item Determine bandwidth, number of antennas and layers per cell and compute the fronthaul rate for candidate splits (\cref{eq:cpri}).
  \item Compare with the available transport (fibre, microwave) and its latency; lower splits need sub-millisecond fronthaul latency.
  \item Weigh the benefits of centralisation (coordination, pooling, simpler sites) against transport cost.
  \item Choose the split per site type (e.g.\ 7.2 for urban massive MIMO, option 2 for remote sites with limited transport).
\end{steps}
\end{procedurebox}

% ---------------------------------------------------------------------
\section{Worked examples}

\begin{examplebox}{Overhead for small packets}\label{ex:oh}
\textbf{Problem.} IoT sensors send 20-byte payloads in IPv6/UDP without header compression. Compute the efficiency on the radio and with ROHC compressing
IPv6/UDP to 4 bytes.

\textbf{Solution.} Without ROHC: $L = 20 + 48 = 68$, data efficiency $20/(68+10) = 25.6\,\%$. With ROHC: $L = 24$, $20/34 = 58.8\,\%$. Header compression
matters most for small packets.
\end{examplebox}

\begin{examplebox}{Fronthaul for massive MIMO}\label{ex:fh}
\textbf{Problem.} A 64-antenna RU with 100\,MHz serves 16 layers. Compare option 8 and option 7.2 fronthaul (assumptions of \cref{sec:fh}).

\textbf{Solution.} Option 8: 260\,Gbit/s. Option 7.2: $273\times12\times28\,000\times16\times18 = 26.4$\,Gbit/s -- ten times less, feasible over a 25 or 50\,Gbit/s Ethernet link.
\end{examplebox}

\begin{examplebox}{Duplication budget}\label{ex:dup}
\textbf{Problem.} Each leg has BLER $10^{-2}$ per transmission and can make 2 transmissions within the budget. What residual loss does a single leg reach, and
what does duplication reach?

\textbf{Solution.} Single leg: $(10^{-2})^2 = 10^{-4}$. With duplication over two independent legs: $10^{-4}\times10^{-4} = 10^{-8}$ (\cref{eq:dup}) -- beyond
the $10^{-5}$ URLLC targets, leaving margin for correlation.
\end{examplebox}

% ---------------------------------------------------------------------
\section{Technology comparison}
\begin{table}[htbp]
\centering\small
\caption{User plane in LTE/EPC and in NR/5GC.}\label{tab:cmp}
\rowcolors{2}{acnmist}{white}
\begin{tabularx}{\linewidth}{@{}l Y Y@{}}
\toprule
\rowcolor{white} & LTE / EPC & NR / 5GC \\
\midrule
Radio layers       & PDCP, RLC (with reordering), MAC & SDAP, PDCP, RLC (no reordering), MAC \\
RAN split          & Monolithic eNB (CPRI to RRH)     & CU/DU (F1), CU-CP/CU-UP (E1), lower-layer split 7.2 \\
Core user plane    & S-GW + P-GW (CUPS optional)      & UPF chains (I-UPF, UL CL, PSA) controlled by PFCP \\
QoS on N3/S1-U     & Tunnel per bearer                & Tunnel per session, QFI per packet \\
Reliability        & HARQ, ARQ                        & + PDCP duplication, multi-connectivity \\
SN length          & PDCP up to 18 bits (later)       & PDCP 12/18, RLC AM 12/18 bits \\
\bottomrule
\end{tabularx}
\end{table}

\paragraph{Programmable all the way} From P4-programmed UPFs to virtualised CU/DU, the 5G user plane is software-defined end to end -- which makes
its performance a matter of engineering budgets (cycles per packet, microseconds per hop) rather than of fixed boxes.

% ---------------------------------------------------------------------
\section{Key takeaways}
\begin{takeawaybox}
\begin{itemize}
  \item User data: UE (SDAP/PDCP/RLC/MAC/PHY) $\to$ DU (RLC/MAC/PHY) $\to$ F1-U $\to$ CU-UP (PDCP/SDAP) $\to$ N3 $\to$ UPF(s) $\to$ N6.
  \item Higher-layer split (option 2) needs about the user rate; option 8 scales with antennas ($2f_s b N_{\mathrm{ant}} c$); 7.2 with layers.
  \item UPF processing: PDR match $\to$ QER $\to$ URR $\to$ FAR, programmed by the SMF over N4.
  \item Latency $= \sum$ components; edge UPFs remove milliseconds of transport.
  \item PDCP duplication gives $P_{\mathrm{loss}} = p_1p_2$; long SNs keep the window $2^{s-1}/P$ above reordering delays.
\end{itemize}
\end{takeawaybox}

% ---------------------------------------------------------------------
\section{Exercises}

\begin{exercisebox}{1}{Which layer?}\label{exr:layer}
Which layer performs: (a) ciphering; (b) segmentation; (c) QoS flow to DRB mapping; (d) HARQ; (e) duplicate removal; (f) ARQ?
\end{exercisebox}

\begin{exercisebox}{1}{Why split?}\label{exr:why}
Give two advantages and two disadvantages of centralising the CU for many DUs.
\end{exercisebox}

\begin{exercisebox}{2}{Fronthaul}\label{exr:fh}
A 40\,MHz carrier at 30\,kHz is sampled at 61.44\,Msample/s. Compute the option-8 fronthaul for 8 antennas with 16-bit samples and 66/64 coding. Can it run over
a 25\,Gbit/s link?
\end{exercisebox}

\begin{exercisebox}{2}{SN window}\label{exr:sn}
A bearer carries 5\,Gbit/s of 1200-byte packets. Compute the window covered by 12-bit and 18-bit PDCP SNs. Which is needed if the reordering timer can reach 100\,ms?
\end{exercisebox}

\begin{exercisebox}{3}{Placing UPFs}\label{exr:place}
An operator has one central data centre (800\,km from the farthest cells) and can build regional sites at most 150\,km from any cell. A service needs 5\,ms one-way
latency with radio 2.5\,ms, gNB 0.5\,ms, UPF and server 0.4\,ms. Which placement works, using 5\,\textmu s/km for transport? Discuss the cost.
\end{exercisebox}

\begin{exercisebox}{3}{Correlated duplication}\label{exr:corr}
Two legs each lose packets with probability $10^{-3}$, but with probability $q$ they are blocked together (common cause), in which case both lose. Model the loss as
$P = q + (1-q)p^2$ with $p = 10^{-3}$ and find the largest $q$ for which $P \le 10^{-5}$. What does this imply for deployment?
\end{exercisebox}

\subsection*{Solutions}
\addcontentsline{toc}{subsection}{Solutions}

\begin{solutionbox}{\ref{exr:layer}}
(a) PDCP; (b) RLC; (c) SDAP; (d) MAC; (e) PDCP; (f) RLC (acknowledged mode).
\end{solutionbox}

\begin{solutionbox}{\ref{exr:why}}
Advantages: pooling of PDCP/RRC processing (statistical multiplexing), easier coordination and mobility between cells of the same CU, simpler and cheaper cell sites,
software upgrades in one place. Disadvantages: transport between CU and DU needed (bandwidth and latency), larger failure domain (one CU outage affects many cells), and
latency-sensitive functions cannot be centralised beyond a distance.
\end{solutionbox}

\begin{solutionbox}{\ref{exr:fh}}
$2\times61.44\times10^6\times16\times8\times66/64 = 16.2$\,Gbit/s: fits a 25\,Gbit/s link for one carrier, but not two such carriers.
\end{solutionbox}

\begin{solutionbox}{\ref{exr:sn}}
$P = 5\times10^9/9600 = 5.2\times10^5$\,packets/s. 12-bit: $2048/P = 3.9$\,ms; 18-bit: $131\,072/P = 0.25$\,s. With a 100\,ms reordering timer, only the 18-bit SN is safe.
\end{solutionbox}

\begin{solutionbox}{\ref{exr:place}}
Fixed part $2.5 + 0.5 + 0.4 = 3.4$\,ms, leaving 1.6\,ms of transport: $D \le 1.6\ \text{ms}/5\ \mu\text{s/km} = 320$\,km. The central site (800\,km, 4\,ms) fails; regional sites
(150\,km, 0.75\,ms) work. Cost: more UPF sites (hardware, operations), regional peering or caching to serve content locally -- to be justified by the revenue of
latency-sensitive services.
\end{solutionbox}

\begin{solutionbox}{\ref{exr:corr}}
$q + (1-q)10^{-6} \le 10^{-5} \Rightarrow q \le 9\times10^{-6}$ (approximately). Common-cause blocking must be extremely rare: the legs should come from different sites and
directions (dual connectivity), not just different carriers of the same antenna.
\end{solutionbox}

% ---------------------------------------------------------------------
\section{Further reading}
The NR user-plane architecture is described in TS~38.300 \cite{ts38300} and the split RAN in TS~38.401 \cite{ts38401}; PFCP in TS~29.244 \cite{ts29244}; GTP-U in TS~29.281
\cite{ts29281}; UPF functions and topologies in TS~23.501 \cite{ts23501}. Dahlman et al.\ \cite{dahlman2020nr} explain the NR layer-2 protocols; Bosshart et al.\
\cite{bosshart2014p4} introduce the P4 language used in programmable UPFs.

\printbibliography[heading=subbibliography,title={\ACNstrRefs}]
\end{document}
